% Source-aligned English for printed p.136.
\SourcePage{136}

\LayerTarget{tohoku:thm:1.10.1:p2}
\noindent from $C$ to $C$, together with a natural transformation $f$ from
the identity functor to $M$, such that for every $A\in C$, $M(A)$ is injective
and $f(A)$ is a monomorphism from $A$ to $M(A)$. Since the proof is
essentially known, we shall sketch only the main points.

\LayerTarget{tohoku:lemma:1.10.1:1}
\noindent\textsc{Lemma 1.} --- \textit{If $C$ satisfies axiom AB 5), then an
object $M\in C$ is injective if and only if, for every subobject $V$ of the
generator $U$ and every morphism $v$ from $V$ to $M$, $v$ extends to a
morphism from $U$ to $M$.}

\LayerTarget{tohoku:lemma:1.10.1:1:proof}
It suffices to prove that the condition is sufficient. Let $u$ be a morphism
to $M$ from a subobject $B$ of an object $A\in C$; we show that it extends to
a morphism from $A$ to $M$. Let $P$ be the set of extensions of $u$ to
subobjects of $A$ containing $B$. This is indeed a set, since the existence of
a generator implies that the subobjects of $A$ form a set. Order $P$ by the
extension relation. By the second formulation of axiom AB 5) (cf. 1.5), this
set is inductive. It therefore has a maximal element. We are thus reduced to
the case in which $u$ itself is maximal, and must prove that $B=A$.

Suppose $B\neq A$. We prove that $u$ extends to some $B'\neq B$. Choose a
morphism $j$ from $U$ to $A$ such that $j(U)\not\subset B$, and put
$B'=j(U)+B$, so $B'\neq B$. Let $V=j^{-1}(B)$ be the inverse image of $B$
under $j$, and let $j':V\longrightarrow B$ be the morphism induced by $j$.
Consider the sequence

\[
 V\mathrel{\mathop{\longrightarrow}^{\varphi'}}U\times B
 \mathrel{\mathop{\longrightarrow}^{\varphi}}B'\longrightarrow0,
\]

\noindent where $\varphi'$ has components the identity map from $V$ to $U$
and $-j'$, while $\varphi$ has components $j$ and the identity map from $B$
to $B'$. This sequence is immediately seen to be exact. Thus, to define a
morphism $v$ from $B'$ to $M$, it suffices to define a morphism $w$ from
$U\times B$ to $M$ such that $w\varphi'=0$. Let $k$ be an extension to $U$ of
$uj':V\longrightarrow M$, and take for $w$ the morphism from $U\times B$ to
$M$ with components $k$ and $u$. Then $w\varphi'=0$, and the morphism
$v:B'\longrightarrow M$ defined by $w$ extends $u$. This proves Lemma 1.

\LayerTarget{tohoku:thm:1.10.1:construction:M1}
Let $A\in C$, and let $I(A)$ be the set of all morphisms $u_i$ from
subobjects $V_i$ of $U$ to $A$. Consider the morphism

\[
 \bigoplus V_i\longrightarrow A\times U^{(I(A))}
\]

\noindent whose restriction to the factor $V_i$ has components
$-u_i:V_i\longrightarrow A$ and the canonical injection of $V_i$ into the
$i$-th factor of the direct sum $U^{(I(A))}$. Let $M_1(A)$ be its cokernel,
and let $f(A):A\longrightarrow M_1(A)$ be the morphism induced by the
canonical epimorphism from $A\times U^{(I(A))}$ onto its quotient. Then
$f(A)$ is a monomorphism; this follows easily from the fact that the canonical
morphism $\bigoplus V_i\longrightarrow U^{(I(A))}$ is a monomorphism by
AB 4). Moreover, every morphism $u_i:V_i\longrightarrow A$ ``extends'' to a
morphism $U\longrightarrow M_1(A)$, namely the morphism induced on the $i$-th
factor of $U^{(I(A))}$ by the canonical epimorphism from
$A\times U^{(I(A))}$ onto its quotient $M_1(A)$.

\LayerTarget{tohoku:thm:1.10.1:construction:transfinite:p1}
By transfinite recursion, for every ordinal $i$ define the object $M_i(A)$,
and for every two ordinals $i\leq j$ an injective morphism
$M_i(A)\longrightarrow M_j(A)$, in such a way that the $M_i(A)$ for $i<i_0$,
where $i_0$ is a fixed ordinal, form an inductive system. For $i=0$ take
$M_0(A)=A$; for $i=1$, $M_1(A)$ and $M_0(A)\longrightarrow M_1(A)$ have
already been defined. If the construction has been made for the ordinals
$<i$, and if $i$ is of the form $j+1$,

\TohokuContinuation{137}
